% s-377-p377-m002; original PDF 377
\PoemStart{s-377-p377-m002}{Die blaue Blume.}
\PoemSubtitle{(Nietzsche an Novalis.)}
\begin{BookVerse}
\VerseLine{0}{Und dir zerfloss die Welt im schimmernd Breiten,}
\VerseLine{1}{In einer mystisch feinen Liebesfluth.}
\VerseLine{0}{Dahin den Strom des Wollens abzuleiten,}
\VerseLine{1}{Darauf verstehst du zarter Geist dich gut.}
\Stanza
\VerseLine{0}{Was musst so fühlsam du Natur besaiten,}
\VerseLine{1}{Dass schon dem Gras die Sichel wehethut?}
\VerseLine{1}{Kaum ritzt der Sporn, so lässt das Rösslein Blut:}
\VerseLine{0}{Wer möchte da ein Ende Wegs noch reiten?}
\Stanza
\VerseLine{0}{Verfallsgewächse zogst du an Spalieren,}
\VerseLine{0}{Die fröstelnd in die Himmels-Heimath stieren}
\VerseLine{0}{Und träumend, eh’ sie noch gelebt, erfrieren.}
\Stanza
\VerseLine{0}{Nicht zu erschlaffen deinen Schlummerflöten,}
\VerseLine{0}{Die schmeichelnd Geist und Kraft und Mannheit tödten,}
\VerseLine{0}{Ist schon allein ein Übermensch von Nöthen.}
\end{BookVerse}
\PoemEnd
